% !PRINTED 114
\weberpage{114}
%<<CONT

\begin{equation}
 \varepsilon_k=\cos\frac{2k\pi}{n}+i\sin\frac{2k\pi}{n}
\tag{11}
\end{equation}
multiplicirt. Die Grössen $\varepsilon_k$ sind Wurzeln der Gleichung
\begin{equation}
 x^n=1,
\tag{12}
\end{equation}
und heissen die $n$ten \gesperrt{Einheitswurzeln}. Es sind $n$ und nur $n$ verschiedene Werthe, von denen bei ungeradem $n$ nur einer (für $k=0$), bei geradem $n$ zwei (für $k=0$ und $k=\frac n2$) reell sind. Die geometrischen Bilder der $n$ten Einheitswurzeln sind die Eckpunkte eines dem Kreise mit dem Radius $1$ einbeschriebenen regelmässigen $n$-Ecks. Ihre algebraische Bestimmung ist Gegenstand der \gesperrt{Kreistheilungslehre}. Bezeichnen wir den Werth von $\varepsilon_k$ für $k=1$ mit $\varepsilon$, setzen also
\[
 \varepsilon=\cos\frac{2\pi}{n}+i\sin\frac{2\pi}{n},
\]
so ist nach dem \gesperrt{Moivre'schen} Satze
\[
 \varepsilon_k=\varepsilon^k,
\]
so dass alle $n$ten Einheitswurzeln als Potenzen von einer unter ihnen dargestellt sind.


\begin{center}
{\large \S.\ 34.}

\vspace{0.5em}
{\large\gesperrt{Befreiung einer Gleichung vom zweiten Gliede.}}
\end{center}

Die Gewissheit der Existenz der Wurzeln reiner Gleichungen setzt uns in den Stand, die Wurzelexistenz bei den Gleichungen zweiten, dritten und vierten Grades oder, wie man auch sagt, bei den quadratischen, cubischen und biquadratischen Gleichungen, nachzuweisen, indem wir die Bestimmung ihrer Wurzeln auf die Lösung reiner Gleichungen zurückführen. Wir werden auf diese Frage später von verschiedenen allgemeineren Standpunkten zurückkommen, besprechen aber hier in der Kürze die älteren Methoden der Auflösung, die, wenn sie auch wenig Einblick in den allgemeinen Zusammenhang dieser Fragen gewähren, doch in der Anwendung sehr einfach sind. Sie erwecken den Schein, als ob es sich um eine der Verallgemeinerung auf höhere Gleichungen fähige Methode handle, was aber nicht der Fall ist. Zunächst folgende allgemeine Bemerkung.

Nehmen wir der Einfachheit halber in der Function
%CONT>>
